\documentclass[11pt]{article}

\usepackage{amsmath, amssymb, amsthm}
\usepackage{geometry}
\usepackage{hyperref}
\usepackage{microtype}
\usepackage{booktabs}

\geometry{margin=1.25in}

\newtheorem{proposition}{Proposition}
\newtheorem{corollary}{Corollary}[proposition]

\title{The Tarzan Scan: Geometry and Periodicity Analysis}
\author{David G.\ Grier}
\date{}

\begin{document}
\maketitle

\begin{abstract}
A Tarzan scan is a polargraph scan pattern in which each move engages
exactly one motor. The payload swings on a circular arc of constant
belt length, tracing the natural coordinate lines of the polargraph.
We derive a closed-form expression for the one-cycle map
$T\colon x_0 \mapsto x_1$ that advances the starting position on
the top edge of the scan rectangle after one complete four-arc cycle.
The map is governed by a single scalar parameter $B$ that depends
only on the scan geometry. When $B = 0$ the map is the identity and
every orbit is period-1; when $B \ne 0$ orbits are generically
aperiodic. We give explicit conditions on the geometry and on $x_0$
that guarantee aperiodic scans.
\end{abstract}

% ------------------------------------------------------------------
\section{Geometry}
% ------------------------------------------------------------------

A polargraph has two pulleys separated by a distance $\ell$,
positioned symmetrically about the origin:
\begin{equation}
  L = \bigl(-h,\,0\bigr), \qquad R = \bigl(h,\,0\bigr),
  \qquad h = \tfrac{\ell}{2}.
\end{equation}
The scan rectangle is specified by a horizontal offset $d_x$, a
vertical offset $d_y$ below the home position $y_0$, a width $w$,
and a height $H$:
\begin{align}
  x_{\mathrm{left}}  &= d_x - w/2, &
  x_{\mathrm{right}} &= d_x + w/2, \\
  y_{\mathrm{top}}   &= y_0 + d_y, &
  y_{\mathrm{bot}}   &= y_0 + d_y + H.
\end{align}
Throughout this document $y$ increases downward (the natural
convention for a hanging polargraph), so
$y_{\mathrm{top}} < y_{\mathrm{bot}}$.

% ------------------------------------------------------------------
\section{One Tarzan Cycle}
% ------------------------------------------------------------------

A single Tarzan cycle starting from $p_0 = (x_0,\, y_{\mathrm{top}})$
consists of four circular arcs, each centred on one pulley:

\begin{center}
\begin{tabular}{cllll}
\toprule
Segment & Centre & Start edge & End edge & Belt fixed \\
\midrule
1 & $R$ & top   & right  & left  \\
2 & $L$ & right & bottom & right \\
3 & $R$ & bottom & left  & left  \\
4 & $L$ & left  & top    & right \\
\bottomrule
\end{tabular}
\end{center}

\noindent
The arc endpoints $p_1, p_2, p_3, p_4$ are computed by intersecting
the relevant circle with the target boundary.

\paragraph{Segment 1 ($R$, top~$\to$ right).}
The belt length from $p_0$ to~$R$ is
$s_R = \lVert p_0 - R \rVert$, so
\begin{equation}
  p_1 = \bigl(x_{\mathrm{right}},\;
              \sqrt{s_R^2 - (x_{\mathrm{right}} - h)^2}\,\bigr).
\end{equation}

\paragraph{Segment 2 ($L$, right~$\to$ bottom).}
The belt length from $p_1$ to~$L$ is $s_L = \lVert p_1 - L \rVert$.
Expanding:
\begin{equation}
  s_L^2
    = (x_{\mathrm{right}} + h)^2 + p_{1y}^2
    = (x_{\mathrm{right}} + h)^2 - (x_{\mathrm{right}} - h)^2
      + s_R^2
    = 4h\,x_{\mathrm{right}} + s_R^2,
\end{equation}
so
\begin{equation}
  p_2 = \bigl(-h + \sqrt{s_L^2 - y_{\mathrm{bot}}^2},\;
              y_{\mathrm{bot}}\bigr).
\end{equation}

\paragraph{Segment 3 ($R$, bottom~$\to$ left).}
\begin{equation}
  p_3 = \bigl(x_{\mathrm{left}},\;
              \sqrt{s_R'^2 - (x_{\mathrm{left}} - h)^2}\,\bigr),
  \qquad
  s_R' = \lVert p_2 - R \rVert.
\end{equation}

\paragraph{Segment 4 ($L$, left~$\to$ top).}
\begin{equation}
  p_4 = \bigl(-h + \sqrt{s_L'^2 - y_{\mathrm{top}}^2},\;
              y_{\mathrm{top}}\bigr),
  \qquad
  s_L' = \lVert p_3 - L \rVert.
\end{equation}

The Tarzan map is $T(x_0) = p_{4x}$.

% ------------------------------------------------------------------
\section{Closed-Form Expression for $T$}
% ------------------------------------------------------------------

Introduce $v = x_0 - h$ and the parameter
\begin{equation}\label{eq:B}
  \boxed{B = 4h\,x_{\mathrm{right}} + y_{\mathrm{top}}^2
             - y_{\mathrm{bot}}^2.}
\end{equation}
We also define the partner parameter
\begin{equation}\label{eq:E}
  E = 4h\,x_{\mathrm{left}} + y_{\mathrm{bot}}^2 - y_{\mathrm{top}}^2,
\end{equation}
which satisfies $B + E = 8h\,d_x$.
For a symmetric scan ($d_x = 0$) we have $E = -B$.

\begin{proposition}\label{prop:map}
The Tarzan map is
\begin{equation}\label{eq:T}
  T(x_0)
  = -h + \sqrt{\bigl(\sqrt{(x_0 - h)^2 + B\,} - 2h\bigr)^2 + E\,}.
\end{equation}
\end{proposition}

\begin{proof}
We track the squared belt lengths through each segment,
writing $v = x_0 - h$.

\medskip\noindent
\textbf{After segment~2.}
$s_R^2 = v^2 + y_{\mathrm{top}}^2$, and from the calculation above,
\begin{equation}
  s_L^2 = 4h\,x_{\mathrm{right}} + s_R^2
         = v^2 + B + y_{\mathrm{bot}}^2,
\end{equation}
so $p_{2x} = -h + \sqrt{v^2 + B}$.

\medskip\noindent
\textbf{After segment~3.}
Let $f = \sqrt{v^2 + B}$, so $p_{2x} = f - h$. Then
\begin{align}
  s_R'^2
    &= (p_{2x} - h)^2 + y_{\mathrm{bot}}^2
     = (f - 2h)^2 + y_{\mathrm{bot}}^2 \notag\\
    &= v^2 + B - 4h f + 4h^2 + y_{\mathrm{bot}}^2.
\end{align}
Expanding $s_R'^2 - (x_{\mathrm{left}} - h)^2$ and collecting terms gives
$p_{3y} = \sqrt{v^2 + D - 4hf}$, where
$D = wh + y_{\mathrm{top}}^2 + 3h^2 - w^2/4$.

\medskip\noindent
\textbf{After segment~4.}
\begin{align}
  s_L'^2
    &= (x_{\mathrm{left}} + h)^2 + p_{3y}^2 \notag\\
    &= (x_{\mathrm{left}} + h)^2 + v^2 + D - 4hf.
\end{align}
The combination $(x_{\mathrm{left}} + h)^2 + D$ simplifies to
\begin{align}
  (x_{\mathrm{left}} + h)^2 + D
    &= 4h(x_{\mathrm{left}} + x_{\mathrm{right}}) + y_{\mathrm{top}}^2
       + 4h^2 \notag \\
    &= 8h\,d_x + y_{\mathrm{top}}^2 + 4h^2,
\end{align}
so
\begin{equation}
  s_L'^2 - y_{\mathrm{top}}^2
    = v^2 + 8h\,d_x + 4h^2 - 4hf
    = (f - 2h)^2 + (8h\,d_x - B) + B - B
\end{equation}
which is $(f - 2h)^2 + E$ after substituting $E = 8h\,d_x - B$.
Therefore
\begin{equation}
  p_{4x}
    = -h + \sqrt{(f - 2h)^2 + E\,}
    = -h + \sqrt{\bigl(\sqrt{v^2 + B\,} - 2h\bigr)^2 + E\,},
\end{equation}
which is equation~\eqref{eq:T} with $v = x_0 - h$.
\end{proof}

% ------------------------------------------------------------------
\section{The Degenerate Case $B = 0$}
% ------------------------------------------------------------------

\begin{proposition}
When $B = 0$ and $d_x = 0$ (so $E = 0$), the map $T$ is the identity:
$T(x_0) = x_0$ for all $x_0 \in [x_{\mathrm{left}}, x_{\mathrm{right}}]$.
\end{proposition}

\begin{proof}
With $B = E = 0$, equation~\eqref{eq:T} becomes
$T(x_0) = -h + |{\,}|x_0 - h| - 2h|$.
For $x_0 \in [x_{\mathrm{left}}, x_{\mathrm{right}}] \subset [-h, h]$
we have $x_0 - h < 0$, so $|x_0 - h| = h - x_0$.
Then $|h - x_0 - 2h| = |{-h - x_0}| = h + x_0$, and
$T(x_0) = -h + (h + x_0) = x_0$.
\end{proof}

The condition $B = 0$ is
\begin{equation}\label{eq:degen}
  \ell w = H\,(y_{\mathrm{top}} + y_{\mathrm{bot}}),
\end{equation}
which relates the motor separation~$\ell$, the scan width~$w$,
the scan height~$H$, and the vertical position of the scan
rectangle. The default parameters of the software's
\texttt{FakePolargraph} ($\ell = 1$\,m, $w = 0.6$\,m,
$H = 0.6$\,m, $y_{\mathrm{top}} = 0.2$\,m,
$y_{\mathrm{bot}} = 0.8$\,m) satisfy this condition exactly, making
every simulated Tarzan scan degenerate.

% ------------------------------------------------------------------
\section{Fixed Points and Periodicity}
% ------------------------------------------------------------------

\begin{proposition}\label{prop:fp}
Suppose $B \ne 0$.
\begin{enumerate}
  \item If $d_x = 0$ then $T$ has no fixed points in
        $[x_{\mathrm{left}}, x_{\mathrm{right}}]$.
  \item If $d_x \ne 0$ then $T$ has exactly one fixed point:
        \begin{equation}\label{eq:fp}
          x_0^* = h + d_x - \frac{B}{4\,d_x}.
        \end{equation}
\end{enumerate}
\end{proposition}

\begin{proof}
Setting $T(x_0) = x_0$ in equation~\eqref{eq:T} and squaring twice
reduces to the linear equation
\begin{equation}
  \sqrt{(x_0 - h)^2 + B\,} = h - x_0 + \frac{B + E}{4h} = h - x_0 + 2d_x.
\end{equation}
Squaring and simplifying using $v = x_0 - h$:
\begin{equation}
  v^2 + B = (2d_x - v)^2 = 4d_x^2 - 4d_x v + v^2,
\end{equation}
so $B = 4d_x(d_x - v) = 4d_x(d_x - x_0 + h)$, giving
$x_0 = h + d_x - B/(4d_x)$ when $d_x \ne 0$.
When $d_x = 0$ this equation becomes $B = 0$, a contradiction.
\end{proof}

\begin{corollary}
For a symmetric scan ($d_x = 0$) with $B \ne 0$, no orbit of~$T$
is period-1.  Because~$T$ is real-analytic, the set of periodic
orbits of any fixed period~$n$ is finite; the union over all~$n$
is countable and therefore has Lebesgue measure zero.  Consequently,
almost every $x_0$ yields an aperiodic scan.
\end{corollary}

% ------------------------------------------------------------------
\section{Practical Criteria for Aperiodic Scans}
% ------------------------------------------------------------------

The following procedure guarantees an aperiodic Tarzan scan.

\begin{enumerate}
  \item \textbf{Break the degeneracy.}
        Compute $B$ using equation~\eqref{eq:B}.
        If $B \approx 0$ (condition~\eqref{eq:degen}), adjust $d_y$
        or~$H$ until $B \ne 0$.  Increasing $d_y$ by a few percent
        of~$H$ is usually sufficient.

  \item \textbf{Use a symmetric scan ($d_x = 0$) if possible.}
        When $d_x = 0$ and $B \ne 0$, Proposition~\ref{prop:fp}
        guarantees no fixed points.  Any choice of $x_0$ in
        $[x_{\mathrm{left}}, x_{\mathrm{right}}]$ will produce an
        aperiodic scan (almost surely).

  \item \textbf{Avoid the fixed point when $d_x \ne 0$.}
        If $d_x \ne 0$, compute $x_0^*$ from
        equation~\eqref{eq:fp} and choose $x_0 \ne x_0^*$.

  \item \textbf{Choose $x_0$ generically.}
        Period-$n$ orbits ($n \ge 2$) form a measure-zero set;
        in practice any $x_0$ not expressible as a simple algebraic
        function of the scan parameters will be aperiodic.
        Starting near $x_{\mathrm{left}} + \epsilon$ for a small
        irrational $\epsilon$ (e.g.\ $\epsilon = \Delta x / \phi$
        where $\phi$ is the golden ratio and $\Delta x$ is the
        step size) is a robust choice.
\end{enumerate}

% ------------------------------------------------------------------
\section{Summary of Parameters}
% ------------------------------------------------------------------

\begin{center}
\begin{tabular}{lll}
\toprule
Symbol & Definition & Role \\
\midrule
$h$ & $\ell/2$ & half motor separation \\
$B$ & $4h\,x_{\mathrm{right}} + y_{\mathrm{top}}^2 - y_{\mathrm{bot}}^2$
    & controls map non-triviality \\
$E$ & $4h\,x_{\mathrm{left}} + y_{\mathrm{bot}}^2 - y_{\mathrm{top}}^2$
    & partner parameter; $B + E = 8h\,d_x$ \\
$v$ & $x_0 - h$ & signed offset from right pulley \\
$T(x_0)$ & $-h + \sqrt{(\sqrt{v^2+B}-2h)^2+E}$
          & one-cycle advance map \\
$x_0^*$ & $h + d_x - B/(4d_x)$ & fixed point (exists iff $d_x \ne 0$, $B \ne 0$) \\
\bottomrule
\end{tabular}
\end{center}

\end{document}
